A general obstacle may have unbounded complex-Hessian coefficients even when its metric Laplacian is bounded. The following estimate nevertheless makes all finite products of these coefficients integrable.

\begin{proposition}\label{prop:compact-metric-laplacian-lp}
Let $1<p<\infty$ and let $f\in L^p(X)$ have $Lf\in L^p(X)$ distributionally. Then $f\in W^{2,p}(X)$ and
\begin{equation}\label{eq:compact-metric-laplacian-lp}
 \|f\|_{W^{2,p}(X)}\leq C_{X,\omega,p}
       \bigl(\|Lf\|_{L^p(X)}+\|f\|_{L^p(X)}\bigr).
\end{equation}
In particular a continuous function with bounded distributional metric Laplacian belongs to $W^{2,p}(X)$ for every finite $p\geq1$ and to $C^{1,\alpha}(X)$ for every $0<\alpha<1$.
\end{proposition}
\begin{proof}
We first prove the estimate for smooth real functions. In real coordinates of dimension $d=2n$, write
\[
 L=-\sum_{i,j}a^{ij}(x)\partial_i\partial_j
               +\sum_i b^i(x)\partial_i,
\]
where the symmetric matrix $a$ is positive definite and all coefficients are smooth. On the compact chart supports its eigenvalues have fixed positive upper and lower bounds. For the frozen operator $L_0=-\sum a^{ij}(x_0)\partial_i\partial_j$, a linear change of coordinates reduces the principal form to the Euclidean one. Equation \eqref{eq:euclidean-hessian-lp} gives
\[
 \|D^2v\|_p\leq C\|L_0v\|_p.
\]
The constant is uniform when the eigenvalues of $a(x_0)$ stay in the fixed bounds. On a sufficiently small ball, $\sup|a(x)-a(x_0)|\leq\delta$, and hence
\[
 \|D^2v\|_p\leq C\bigl(\|Lv\|_p+
           \delta\|D^2v\|_p+\|\nabla v\|_p\bigr)
\]
for $v$ compactly supported there. Choose $\delta$ to absorb the Hessian term.

Take a finite smooth partition $\chi_\nu$ subordinate to these balls and apply the estimate to $v_\nu=\chi_\nu f$, extended by zero in its chart. The commutators $[L,\chi_\nu]f$ contain only first derivatives of $f$. Summing gives
\[
 S_2:=\sum_\nu\|D^2v_\nu\|_p
 \leq C\bigl(\|Lf\|_p+\|\nabla f\|_p+\|f\|_p\bigr).
\]
All chart and metric norms here are equivalent on the fixed finite cover. Viewing the localized functions again on $X$, we have $f=\sum_\nu v_\nu$; its gradient norm is at most a fixed multiple of $\sum_\nu\|\nabla v_\nu\|_p$. By \cref{lem:gradient-lp-interpolation}, this is at most $C\varepsilon S_2+C_\varepsilon\|f\|_p$. Absorption gives \eqref{eq:compact-metric-laplacian-lp} for smooth functions, including their first derivatives.

Let $P=L+c$ with the fixed $c>0$ used in the proof of \cref{lem:compact-elliptic-foundations}. For smooth real $b$, that proof gives a smooth real solution $u=P^{-1}b$. Testing the equation by
\[
 \phi_\varepsilon(u)=u(u^2+\varepsilon)^{(p-2)/2},
 \qquad \varepsilon>0,
\]
is legitimate. Its derivative is
\[
 \phi_\varepsilon'(u)=
 ((p-1)u^2+\varepsilon)(u^2+\varepsilon)^{(p-4)/2}\geq0.
\]
Stokes therefore gives $c\int_Xu\phi_\varepsilon(u)\,\mathrm dV\leq\int_Xb\phi_\varepsilon(u)\,\mathrm dV$. Smoothness and compactness permit dominated convergence as $\varepsilon\downarrow0$, also for $1<p<2$, where $|\phi_\varepsilon(u)|\leq|u|^{p-1}$. H\"older gives
\[
 \|P^{-1}b\|_p\leq c^{-1}\|b\|_p.
\]
The smooth estimate already proved, applied with $Lu=b-cu$, now yields $\|P^{-1}b\|_{W^{2,p}}\leq C\|b\|_p$.

For the given real $f$ put $b=Lf+cf\in L^p$, and choose smooth real $b_j\to b$ in $L^p$ by chart cutoffs and mollification. The solutions $u_j=P^{-1}b_j$ are Cauchy in $W^{2,p}$. Completeness follows coefficientwise from completeness of $L^p$, with the weak derivatives identified by testing. Their limit $u$ satisfies $Pu=b$ distributionally. The distributional inverse identities constructed in \cref{lem:compact-elliptic-foundations} give uniqueness for $P$, so $u=f$. This proves the estimate. Apply the real assertion to the real and imaginary parts for complex $f$.

If $f$ is continuous and $Lf$ is bounded, both are in every finite $L^p$ on $X$. The assertion for $p=1$ follows from $W^{2,2}$ and finite volume. Given $0<\alpha<1$, choose $p>d/(1-\alpha)$. Morrey's inequality, \cref{lem:morrey-sobolev-representative}, in each chart gives a $C^{1,1-d/p}$ representative, hence $C^{1,\alpha}$. It equals the original continuous function because the two agree almost everywhere.
\end{proof}
\index{Laplacian!$W^{2,p}$ regularity}

\begin{corollary}\label{cor:metric-laplacian-barrier-determinants}
Let $\theta_1,\ldots,\theta_n$ be closed smooth real $(1,1)$-forms on $X$, and let $f_1,\ldots,f_n$ be continuous real functions with bounded distributional $Lf_j$. Their weak forms have a well-defined wedge
\[
 \bigwedge_{j=1}^n(\theta_j+\ddc f_j)
       =h\,\mathrm dV,\qquad h\in L^q(X)\quad(1\leq q<\infty).
\]
This is a finite signed Radon measure. On a coordinate ball choose smooth real $\rho_j$ with $\theta_j=\ddc\rho_j$; there it is $n!\pi^{-n}D(B_1,\ldots,B_n)\,\mathrm d\lambda$, with $B_j=(\partial_a\bar\partial_b(\rho_j+f_j))$. If continuous $\theta_j$-psh functions $u_j$ also have bounded $Lu_j$, then on $E=\bigcap_j\{u_j=f_j\}$ this signed measure agrees with the positive contact measure in \cref{prop:metric-laplacian-regular-contact-density}.
\end{corollary}
\begin{proof}
The proposition puts all weak second derivatives of $f_j$ in every finite $L^p$. Products of $n$ coefficients are therefore in $L^q$ by H\"older with $p=nq$; when $nq=1$, use $p=2$ and finite volume instead. Expansion of the wedges defines their density. Weak differential forms transform by the smooth coordinate rule, so these polynomial expressions agree on chart overlaps almost everywhere. The normalization is the smooth determinant calculation in \cref{prop:bounded-complex-hessian-density}, applied algebraically to the weak matrices. A real $L^1$ density defines a finite signed Radon measure by its positive and negative parts. The contact identity is precisely \cref{prop:metric-laplacian-regular-contact-density}.
\end{proof}
\index{Monge--Amp\`ere measure!weak barrier determinant}
